%!TEX program = xelatex
% !TEX encoding = ISO-8859-1
\def\ano{2019}
\def\semestre{2}
\def\disciplina{\'Algebra 1}
\def\turma{C}

\documentclass[12pt]{exam}

\usepackage{caption}
\usepackage{amssymb}
\usepackage{amsmath,amsfonts,amsthm,amstext}
\usepackage[brazil]{babel}
% \usepackage[latin1]{inputenc}
\usepackage{graphicx}
\graphicspath{{/home/jfreitas/GitHub_Repos/Algebra-1/Pictures/}{D:/OneDrive - unb.br/imagens-latex/}}
\usepackage{enumitem}
\usepackage{multicol}
\usepackage{answers}
\usepackage{tikz,ifthen}
\usetikzlibrary{lindenmayersystems}
\usetikzlibrary[shadings]
\Newassociation{solucao}{Solution}{ans}
\newtheorem{exercicio}{}

\setlength{\topmargin}{-1.0in}
\setlength{\oddsidemargin}{0in}
\setlength{\textheight}{10.1in}
\setlength{\textwidth}{6.5in}
\setlength{\baselineskip}{12mm}

\extraheadheight{0.7in}
\firstpageheadrule
\runningheadrule
\lhead{
        \begin{minipage}[c]{1.7cm}
        \includegraphics[width=1.7cm]{unb.pdf}
        \end{minipage}%
        \hspace{0pt}
        \begin{minipage}[c]{4in}
          {Universidade de Bras{\'\i}lia} --
          {Departamento de Matem{\'a}tica}
\end{minipage}
\vspace*{-0.8cm}
}
% \chead{Universidade de Bras{\'\i}lia - Departamento de Matem{\'a}tica}
% \rhead{}
% \vspace*{-2cm}

\extrafootheight{.5in}
\footrule
\lfoot{\disciplina\ - \semestre$^o$/\ano}
\cfoot{}
\rfoot{P\'agina \thepage\ de \numpages}

\newcounter{exercicios}
\renewcommand{\theexercicios}{\arabic{exercicios}}

\newenvironment{questao}[1]{
\refstepcounter{exercicios}
\ifx&#1&
\else
   \label{#1}
\fi
\noindent\textbf{Exerc{\'\i}cio {\theexercicios}:}
}

\newcommand{\resp}[1]{
\noindent{\bf Exerc{\'\i}cio #1: }}
\newcommand{\nsub}{\varsubsetneq}
\newcommand{\vaz}{\emptyset}
\newcommand{\im}{{\rm Im\,}}
\newcommand{\sub}{\subseteq}
\newcommand{\n}{\mathbb{N}}
\newcommand{\z}{\mathbb{Z}}
\newcommand{\rac}{\mathbb{Q}}
\newcommand{\real}{\mathbb{R}}
\newcommand{\complex}{\mathbb{C}}
\newcommand{\cp}[1]{\mathbb{#1}}
\newcommand{\ch}{\mbox{\textrm{car\,}}\nobreak}
\newcommand{\vesp}[1]{\vspace{ #1  cm}}
\newcommand{\compcent}[1]{\vcenter{\hbox{$#1\circ$}}}
\newcommand{\comp}{\mathbin{\mathchoice
{\compcent\scriptstyle}{\compcent\scriptstyle}
{\compcent\scriptscriptstyle}{\compcent\scriptscriptstyle}}}

\begin{document}

\begin{center}
{\Large\bf \disciplina\ - Turma \turma\ -- \semestre$^{o}$/\ano} \\ \vspace{9pt} {\large\bf
$2^{\underline{o}}$ Teste - M\'odulo 2 - Resolu\c{c}\~ao}\\
\vspace{9pt} Prof. Jos{\'e} Ant{\^o}nio O. Freitas
\end{center}
\hrule

\vspace{.6cm}

\questao Seja $f : A \to B$ uma fun{\c c}{\~a}o e seja $X \sub B$. Mostre que $f^{-1}(X^C) = (f^{-1}(X))^C$.
\noindent\textbf{Solu\c{c}\~ao:} Precisamos mostrar que
\begin{enumerate}
    \item $f^{-1}(X^C) \sub (f^{-1}(X))^C$
    \item $(f^{-1}(X))^C \sub f^{-1}(X^C)$.
\end{enumerate}

Para a primeira inclusão seja $z \in f^{-1}(X^C)$. Então, por definição, $f(z) \in X^C$. Logo, $f(z) \notin X$. Com isso $z \notin f^{-1}(X)$, ou seja, $z \in (f^{-1}(X))^C$.

Agora para a segunda inclusão, seja $t \in (f^{-1}(X))^C$. Assim $t \notin f^{-1}(X)$ e então $f(t) \notin X$. Daí $f(t) \in X^C$, isto é, $t \in f^{-1}(X^C)$.

Portanto, $f^{-1}(X^C) = (f^{-1}(X))^C$.


\vspace{.5cm}

\questao Seja
\[
    M_2(\real) = \left\{\begin{pmatrix}
        x & y\\z & t 
    \end{pmatrix} \mid x, y, z, t \in \real \right\}.
\]
Dados $A = \begin{pmatrix}
        x & y\\z & t 
    \end{pmatrix}$, $B = \begin{pmatrix}
        a & b\\c & d 
    \end{pmatrix} \in M_2(\real)$ definimos
\begin{align*}
    A &= B \mbox{ quando } x = a, y = b, z = c, t = d\\
    A + B &= \begin{pmatrix}
        x & y\\z & t 
    \end{pmatrix} + \begin{pmatrix}
        a & b\\c & d 
    \end{pmatrix} = \begin{pmatrix}
        x + a & y + b\\z + c & t + d 
    \end{pmatrix}\\
    A \cdot B &= \begin{pmatrix}
        x & y\\z & t 
    \end{pmatrix} \cdot \begin{pmatrix}
        a & b\\c & d 
    \end{pmatrix} = \begin{pmatrix}
        xa + yc & xb + yd\\za + tc & zb + td 
    \end{pmatrix}.
\end{align*}

Mostre que $(M_2(\real), + , \cdot)$ \'e um an\'el. Esse anel \'e comutativo? Possui unidade?
\end{document}